\documentclass[12pt]{article}
\usepackage[utf8]{inputenc}
\usepackage[T1]{fontenc}
\usepackage{hyperref}
\title{Togoyo Language Resources for K-12 Curriculum Development}
\author{Compiled by Claude Sonnet 5 (claude-sonnet-5) for Andr\'as Kornai}
\date{}
\begin{document}
\maketitle
\section{Introduction}
Togoyo (ISO 639-3: tgy) is an extinct Ubangian (Sere--Mba) language formerly spoken in Raga County, Western Bahr al Ghazal, South Sudan. Sources give its extinction as the mid-1980s (ASJP: 1985), with about 700 speakers estimated before it died out. Very little material exists, and no native-language educational resources were found in the searches.
\section{Primary Educational Resources}
No primers, textbooks or school programmes were found for Togoyo.
\section{Linguistic and Reference Materials}
\begin{itemize}
\item Wikipedia overview of the Togoyo language, which cites Santandrea's 1953 article on Indri, Togoyo, Feroge, Mangaya and Woro: \url{https://en.wikipedia.org/wiki/Togoyo_language}
\item Wikipedia, Sere languages (classification; Indri--Togoyo): \url{https://en.wikipedia.org/wiki/Sere_languages}
\item MPI numerals database entry for Togoyo (numeral data, speaker estimate): \url{https://lingweb.eva.mpg.de/channumerals/Togoyo.htm}
\item ASJP database entry (word list data): \url{https://asjp.clld.org/languages/TOGOYO}
\end{itemize}
\section{Religious and Community Resources}
None found.
\section{Digital and Audio Resources}
No audio or video found. The ASJP word list is the only lexical data located in the search results.
\section{Curriculum Development Context}
With no living speakers and only a short word list and a 1953 source description, a curriculum would be a heritage-awareness unit rather than a language course. Primary sources (Santandrea 1953) would need to be located separately and checked.
\end{document}